\chapter{Mamba}
\label{ch:mamba}

\section{Why Mamba Was Needed}

S4 and S5 showed that structured state space models can supply long memory with linear
time scaling and fixed-size recurrent state. But they also inherit a major restriction:
their temporal core is linear time invariant. The same kernel is applied regardless of
the input content.

That is a real limitation. In many sequence tasks, the model should not treat every token
or time step the same way. Some inputs should be remembered strongly; others should be
ignored. Some events should reset memory; others should be carried forward. Attention does
this through content-dependent weights. Plain LTI SSMs do not.

Mamba addresses this gap by making the state space model \emph{selective}
\citep{gu2023mamba}. Its key idea is:
\[
  \text{the recurrence parameters depend on the current input.}
\]
This gives the model a form of content-based routing while keeping a linear-time scan-like
computation.

\section{From Fixed Dynamics To Selective Dynamics}

Recall the standard discrete SSM:
\[
  \vx_t = \bar{\mA}\vx_{t-1} + \bar{\mB}\vu_t,
  \qquad
  \vy_t = \mC\vx_t + \mD\vu_t.
\]
In S4 and S5, the matrices are fixed across time. The same $\bar{\mA}$, $\bar{\mB}$,
and $\mC$ are used for every $t$.

Mamba instead uses input-dependent parameters. A simplified conceptual form is
\begin{equation}
  \vx_t = \bar{\mA}_t \vx_{t-1} + \bar{\mB}_t \vu_t,
  \qquad
  \vy_t = \mC_t \vx_t.
  \label{eq:mamba-selective-ssm}
\end{equation}
The matrices at time $t$ are functions of the current input or a projection of it. In the
original Mamba formulation, the continuous-time state matrix itself stays structured, but
the discretization step size and input/output maps are input-dependent. The effect is still
that the discrete-time recurrence changes from one step to the next.

This is the meaning of \emph{selection}. The model selects how to write to memory, how
fast to update, and how to read out, based on the current input.

\section{Why Fixed Convolution Breaks}

For an LTI SSM, we can unroll the recurrence into a single kernel:
\[
  \mK_i = \mC\bar{\mA}^i\bar{\mB}.
\]
That works only because the parameters are fixed across time.

In Mamba, if the recurrence changes with $t$, then the effect of an old input on the
current output depends on the whole chain of intermediate updates:
\[
  \vy_t
  \text{ depends on }
  \bar{\mA}_t \bar{\mA}_{t-1}\cdots \bar{\mA}_{s+1}
  \bar{\mB}_s \vu_s.
\]
There is no single fixed kernel depending only on lag $(t-s)$. The system is no longer
time invariant. So ordinary convolution mode disappears.

This is the key tradeoff:

\[
\begin{array}{lll}
\text{S4/S5} &:& \text{fixed kernel, strong convolution view, no content-based selection}, \\
\text{Mamba} &:& \text{input-dependent recurrence, no fixed kernel, selective memory}.
\end{array}
\]

Mamba gives up one kind of algebraic convenience in exchange for a more flexible temporal
mechanism.

\section{What Is Being Selected?}

The original Mamba paper introduces selective state spaces by making certain SSM
components input-dependent, especially:

\begin{itemize}
  \item the step size $\Delta_t$;
  \item the input matrix or input vector $\mB_t$;
  \item the output readout $\mC_t$.
\end{itemize}

The conceptual effect of each is different.

\paragraph{Selective step size.}
If $\Delta_t$ changes with the input, then the effective decay or timescale of memory can
change at each step. A large step may cause fast update or forgetting. A small step may
preserve state longer.

\paragraph{Selective input write-in.}
If $\mB_t$ changes with the input, the model can decide what kind of information enters the
state and which modes it excites.

\paragraph{Selective readout.}
If $\mC_t$ changes with the input, the model can decide which state components matter for
the current output.

Together, these give the model something LTI SSMs lack:
\[
  \text{input-dependent control over memory flow.}
\]

\section{A Useful Intuition}

Think of S5 as a bank of fixed memory modes. Once the mode spectrum is learned, every time
step is processed through the same memory mechanism. The input changes the state values,
but not the update rule itself.

Mamba instead says: the update rule can change at each step. Some inputs can tell the
model to write aggressively, some to preserve state, some to expose certain modes in the
output, and some to suppress them.

This is not the same as attention, but it serves a related purpose. Attention performs
content-based routing through explicit pairwise interactions. Mamba performs a different
kind of content-based routing through input-conditioned state updates.

\section{The Selective Scan}

At first glance, input-dependent recurrence seems to destroy the efficiency story. If the
matrices change at every step, are we back to a purely sequential RNN?

Mamba's answer is no. Even though there is no fixed convolution kernel, the recurrence
still has an associative structure that can be exploited with a scan algorithm. This is
often called the \emph{selective scan}.

The recurrence at each time step defines an affine map:
\[
  F_t(\vx)=\bar{\mA}_t\vx+\vb_t,
  \qquad
  \vb_t=\bar{\mB}_t\vu_t.
\]
As in Chapter~\ref{ch:s5}, affine-map composition is associative:
\[
  (\mA_2,\vb_2)\circ(\mA_1,\vb_1)
  =
  (\mA_2\mA_1,\mA_2\vb_1+\vb_2).
\]
This means prefix compositions can still be parallelized. The resulting implementation is
not the same as FFT convolution, but it still scales linearly in sequence length and can
be hardware-efficient.

\section{Hardware-Aware Design}

The Mamba paper is not only about mathematical recurrence. It is also about GPU-friendly
implementation. The selective scan is designed so that the recurrence can be evaluated in
a memory-efficient, fused way rather than through slow generic loops.

This is part of why Mamba matters in practice. A model can have attractive theory and
still lose on real hardware if its kernels are badly matched to accelerator execution.
Mamba's engineering contribution is substantial:

\[
  \text{make selective recurrent computation fast enough to compete in practice.}
\]

For learning purposes, the important idea is that algorithm design and model design are
tightly coupled here. Mamba is not ``just S5 with input dependence.'' It is a specific
input-dependent SSM together with a computation strategy that preserves throughput.

\section{The Mamba Block}

A Mamba block wraps the selective SSM inside a neural-network block. A simplified view is:
\[
  \vz
  \rightarrow
  \text{input projection}
  \rightarrow
  \text{split into branches}
  \rightarrow
  \text{local convolution}
  \rightarrow
  \text{selective SSM scan}
  \rightarrow
  \text{gating}
  \rightarrow
  \text{output projection}
  \rightarrow
  \text{residual addition}.
\]

The exact block contains several learned projections. One branch prepares features for the
selective SSM. Another branch provides a multiplicative gate. There is also a short local
convolution before the selective scan.

This local convolution matters. It gives Mamba a small amount of explicit local mixing in
addition to the recurrent state mechanism. The block is therefore not purely ``SSM only.''
It combines:

\begin{itemize}
  \item local short-range convolution;
  \item long-range selective recurrence;
  \item channel projections and gating.
\end{itemize}

\section{Why The Local Convolution Is There}

The selective SSM is good at long-range memory, but local pattern extraction is also
important in sequence data. A short convolution near the input can cheaply capture local
edges, short motifs, or small-scale temporal structure before the state mechanism handles
longer-range effects.

This is a common pattern in sequence architectures. Even transformers often include
positional or local mixing components. Mamba makes this explicit inside the block.

For neural data, this local component is plausible. Utah-array sequences often contain
short-timescale temporal structure that may be easier to capture with local filters before
or alongside longer dynamical memory.

\section{Gating And Feature Control}

The multiplicative gate in a Mamba block plays a role analogous to gating mechanisms in
other architectures: it lets the block modulate how much of the temporally mixed signal is
passed forward.

At a high level:
\[
  \text{output} \approx \text{gate} \odot \text{selective-SSM-output}.
\]
This is not the full formula, but it captures the intent. The gate helps control feature
flow and increases expressive power without abandoning the efficient recurrent structure.

\section{Mamba Versus Attention}

The cleanest comparison is functional rather than ideological.

\paragraph{Attention.}
Attention gives direct pairwise access from one position to another. The weight from $t$ to
$s$ is content-dependent and can vary arbitrarily across pairs. This is very flexible, but
causal self-attention over long sequences is memory-intensive.

\paragraph{Mamba.}
Mamba does not create arbitrary pairwise interactions. It processes the sequence through a
state. The current input influences how the state is updated and read. This is less
flexible than arbitrary attention, but much cheaper in long-sequence scaling and naturally
online.

So Mamba is best thought of as:
\[
  \text{more selective than S4/S5, less pairwise explicit than attention.}
\]

\section{Mamba Versus S4 And S5}

Relative to S4 and S5, Mamba changes the computational center of gravity.

\[
\begin{array}{lll}
\text{S4} &:& \text{structured LTI SSM, convolution-heavy viewpoint}, \\
\text{S5} &:& \text{diagonal LTI SSM, scan-friendly recurrence viewpoint}, \\
\text{Mamba} &:& \text{selective SSM, input-dependent scan viewpoint}.
\end{array}
\]

S4 and S5 are easiest to analyze as kernels or fixed modes. Mamba is easier to analyze as
a learned, gated, input-conditioned recurrent system. Once selectivity is introduced, the
static kernel picture is no longer the right primary mental model.

\section{Why Mamba Can Handle Content Better}

Consider a simple language-style example. If the current token signals the start of a new
entity, the model might want to preserve certain information for a long time. If the token
signals a boundary or reset, it might want to forget. Attention can do this through
content-based weights. Mamba can do it by changing the recurrence parameters.

For time series, the analogue is:

\begin{itemize}
  \item some inputs indicate meaningful events worth remembering;
  \item some inputs indicate nuisance fluctuations worth down-weighting;
  \item some inputs suggest the model should read out one timescale more strongly than
  another.
\end{itemize}

This is the practical meaning of selectivity. The model is not just integrating signal
with a fixed filter. It is deciding how to process the signal based on its content.

\section{Limitations Of Mamba}

Mamba is powerful, but it also brings tradeoffs.

\paragraph{Harder theory than LTI SSMs.}
With fixed kernels, it is easy to inspect modes, impulse responses, and transfer functions.
With input-dependent dynamics, analysis is less clean.

\paragraph{Less direct interpretability.}
A fixed diagonal SSM can be described as a bank of modes. A selective SSM changes its
behavior from one step to the next, which is harder to summarize.

\paragraph{Still not full attention.}
Mamba introduces selectivity, but it does not explicitly compute arbitrary pairwise
token-token interactions. There are tasks where full attention may still be a better fit.

\paragraph{Objective dependence remains.}
As with S5, a strong backbone does not rescue a badly chosen SSL objective. The model may
still learn shortcuts if the pretraining task allows them.

\section{Where Mamba Fits In The Story}

Mamba is best understood as the next step after S4 and S5:

\begin{enumerate}
  \item RNNs gave recurrent state but poor parallelism.
  \item CNNs gave parallel filters but fixed local kernels.
  \item S4 and S5 gave long-range structured recurrent filters with efficient training.
  \item Mamba added content-dependent selectivity while preserving linear-time scan-style
  computation.
\end{enumerate}

That is why Mamba became so influential. It offered a plausible answer to the criticism
that LTI SSMs are too rigid, while keeping the efficient state-space deployment story.

\section{A Brief Note On Mamba-2}

The follow-up work often called Mamba-2 or structured state space duality
\citep{dao2024mamba2} pushes the theory further. It strengthens the connection between
SSMs and attention-like sequence modeling and refines the computational picture.

For a first pass, however, the original Mamba paper is enough. The main conceptual gain is
already there:
\[
  \text{state-space sequence modeling can be selective, not just fixed and linear time invariant.}
\]

\section{What To Remember}

\begin{itemize}
  \item Mamba is a selective state space model.
  \item It introduces input-dependent recurrence parameters, especially in step size and
  input/output mappings.
  \item Because the dynamics depend on the input, there is no single fixed convolution
  kernel.
  \item Mamba therefore relies on selective scan rather than ordinary convolution mode.
  \item The Mamba block combines local convolution, selective recurrence, gating, and
  channel projection.
  \item Mamba sits between LTI SSMs and attention: more content-aware than S4/S5, less
  explicitly pairwise than self-attention.
  \item Its practical importance is as much about efficient implementation as about model
  equations.
\end{itemize}

\section{Exercises}

\begin{enumerate}
  \item Explain why input-dependent $\bar{\mA}_t$ destroys the existence of a fixed kernel
  depending only on lag.
  \item Compare the kind of content dependence used by attention with the kind used by
  Mamba.
  \item Why is associative scan still possible even when the recurrence changes with time?
  \item In your own words, what modeling role is played by the local convolution inside a
  Mamba block?
  \item Why might Mamba be especially appealing for long sequences that still need
  content-sensitive memory updates?
\end{enumerate}

\section{Prepares You For}

The final comparison chapter places Mamba alongside CNNs, RNNs, transformers, S4, and S5,
with special attention to Utah array self-supervised learning.
